\chapter{Design}

\section{General architecture}

The platform is organized into four main components, communicating over HTTP APIs:
\begin{itemize}
\item a \textbf{frontend} built with Next.js, offering three distinct spaces depending on the role (Administrator, Head Technician, Technician);
\item a \textbf{backend} built with FastAPI, responsible for business logic, data persistence (PostgreSQL), and security (authentication, access control);
\item a \textbf{Machine Learning microservice} built with Flask, in charge of serving the predictive models (P1 to P7);
\item a \textbf{RAG service}, in charge of document ingestion and generating conversational responses enriched with real-time ML context.
\end{itemize}

The whole system is orchestrated by Docker, with object storage (MinIO) for documents and attachments.

\begin{figure}[h]
\centering
\begin{tikzpicture}[node distance=1.3cm, box/.style={draw, rectangle, minimum width=3.2cm, minimum height=1cm, align=center}]
\node[box] (front) {Frontend\\Next.js};
\node[box, below=of front] (back) {Backend\\FastAPI};
\node[box, below left=1cm and 0.5cm of back] (ml) {ML Microservice\\Flask};
\node[box, below right=1cm and 0.5cm of back] (rag) {RAG Service};
\node[box, below=2.6cm of back] (db) {PostgreSQL / MinIO};
\draw[->] (front) -- (back);
\draw[->] (back) -- (ml);
\draw[->] (back) -- (rag);
\draw[->] (ml) -- (db);
\draw[->] (rag) -- (db);
\draw[->] (back) -- (db);
\end{tikzpicture}
\caption{General architecture of the platform}
\label{fig:architecture}
\end{figure}

\section{Global business process flow}

Beyond the component-level architecture above, it is useful to see how a single business event travels through the whole system. An intervention request follows one validation gate before becoming a work order, while machine telemetry is analyzed continuously by the ML engine, whose predictions feed directly into that same validation step as supporting evidence for the Head Technician.

\begin{figure}[h]
\centering
\includegraphics[width=0.85\textwidth]{figures/global.png}
\caption{Global flow of an intervention request}
\label{fig:flux-global}
\end{figure}

Rejecting a request (dashed box in Figure~\ref{fig:flux-global}) automatically releases any parts reservation already made for it; approving one creates a linked work order and moves execution forward until a second, final validation closes the diagnostic loop before an Administrator can close the order.

\section{Data model}

The data model is built around four central entities, identified as the most strongly connected during the domain analysis:
\begin{itemize}
\item \textbf{User}: carries the role (Administrator, ChefTech, Technician) and access rights;
\item \textbf{Machine}: represents a piece of production equipment, with its characteristics and history;
\item \textbf{WorkOrder}: a work request associated with one or more machines;
\item \textbf{InterventionOrder}: a concrete intervention carried out by a technician on a machine, linked to a work order.
\end{itemize}

An \textbf{Alert} entity connects these objects to the ML microservice's predictions.

\begin{figure}[h]
\centering
\begin{tikzpicture}[node distance=1.5cm, box/.style={draw, rectangle, minimum width=3cm, minimum height=1cm, align=center}]
\node[box] (user) {User};
\node[box, right=2.5cm of user] (machine) {Machine};
\node[box, below=of user] (ot) {WorkOrder};
\node[box, below=of machine] (oi) {InterventionOrder};
\node[box, below right=1.2cm and -0.5cm of ot] (alerte) {Alert};
\draw (user) -- (ot);
\draw (machine) -- (ot);
\draw (ot) -- (oi);
\draw (machine) -- (oi);
\draw (oi) -- (alerte);
\draw (machine) -- (alerte);
\end{tikzpicture}
\caption{Simplified data model}
\label{fig:modele-donnees}
\end{figure}

\section{Machine Learning pipeline architecture}

The prediction pipeline is structured into specialized models:
\begin{itemize}
\item \textbf{P1}: failure probability;
\item \textbf{P2}: failure type classification;
\item \textbf{P3}: remaining useful life (RUL) estimation, using a hybrid model combining a recurrent network and a boosted decision tree;
\item \textbf{P4}: behavioral anomaly detection, using an ensemble of four combined methods (isolation forest, Z-score, cluster deviation, autoencoder);
\item \textbf{P5}: work order prioritization;
\item \textbf{P6}: maintenance scheduling.
\end{itemize}

These seven models are trained independently and never consume one another's output: each reads the same shared telemetry feature set and returns its own field in the API response, as shown in Figure~\ref{fig:p1-p7}.

\begin{figure}[h]
\centering
\includegraphics[width=\textwidth]{figures/p1p7.png}
\caption{The seven independent P1--P7 prediction models}
\label{fig:p1-p7}
\end{figure}

A second, separate pipeline (referred to as ``Wave 2'') reads each machine's telemetry \emph{history} through four signal-processing models: a physics-informed RUL estimator, a Cox proportional-hazards survival model, a Mahalanobis health index, and a CUSUM anomaly detector smoothed by a Kalman filter. Their outputs are converted into probability assignments over four hypotheses (Healthy, Degrading, Critical, Unknown) and combined through Dempster-Shafer theory fusion --- falling back to the more conservative Yager's rule when the models disagree strongly --- producing the single \emph{unified health score} shown to users, together with an overall verdict. This score, not the P1--P7 outputs, is what downstream readiness and alerting logic consumes, as shown in Figure~\ref{fig:dst-fusion}.

\begin{figure}[h]
\centering
\includegraphics[width=\textwidth]{figures/dst.png}
\caption{Wave 2 signal fusion into the DST health score}
\label{fig:dst-fusion}
\end{figure}

The readiness score blends the unified health score (40\%) with inventory coverage (30\%), procurement risk (20\%), and maintenance recency (10\%). Alerts aggregate signals from both pipelines --- an anomaly flag from P4, a parts shortage from P7, or a model-disagreement flag from the fusion step itself --- and a detected parts shortage is what triggers the purchase draft, as shown in Figure~\ref{fig:downstream}.

\begin{figure}[h]
\centering
\includegraphics[width=0.85\textwidth]{figures/downstream.png}
\caption{From ML outputs to readiness, alerts, and purchase drafts}
\label{fig:downstream}
\end{figure}

A final module, \textbf{P7}, uses failure predictions to anticipate spare parts demand (through survival modeling combined with an intermittent demand forecasting method), and proposes guided --- never automatic --- provisioning: every draft still requires explicit human approval before anything is ordered.

\section{RAG service and ML bridge architecture}

The RAG service ingests technical documentation (maintenance manuals, safety procedures, machine data sheets) and splits it into indexed fragments for semantic search. When a technician asks a question about a specific machine, the backend queries the RAG service (documentation) and the ML microservice (the machine's real-time health status) in parallel, then merges the two contexts before passing them to the language model to generate a response, as shown in Figure~\ref{fig:rag-bridge}. This mechanism ensures the assistant always responds, even if one of the two services is unavailable.

\begin{figure}[h]
\centering
\includegraphics[width=0.85\textwidth]{figures/rag.png}
\caption{RAG service and ML bridge}
\label{fig:rag-bridge}
\end{figure}

\section{Technology choices}

\begin{itemize}
\item \textbf{Next.js}: hybrid rendering, mature React ecosystem, well suited to rich, role-based interfaces;
\item \textbf{FastAPI}: high performance, native schema validation, automatically generated API documentation;
\item \textbf{Flask}: lightweight, well suited to a model-serving microservice;
\item \textbf{PostgreSQL}: transactional robustness, relational queries well suited to the data model;
\item \textbf{Docker}: isolation and reproducibility of development and deployment environments;
\item \textbf{MinIO}: S3-compatible, self-hostable object storage.
\end{itemize}
